\section*{第 \ref{ch:b3:rate-theories} 章——反应速率理论}

\begin{solution}{exo:b3:rate-theories:1}
$\bar v_{\mathrm{rel}} = \sqrt{8 \times 1.381\times10^{-23} \times 298/(\pi \times 18.46 \times 1.661\times10^{-27})} = \qty{585}{m/s}$.
$\sigma\bar v_{\mathrm{rel}}N_A = 0.43\times10^{-18} \times 585 \times 6.022\times10^{23} = \qty{1.51e8}{m^3.mol^{-1}.s^{-1}} =
\qty{1.5e11}{L.mol^{-1}.s^{-1}}$。实测值：$2.07\times10^{-12} \times 6.022\times10^{23}\times10^{-3} =
\qty{1.25e9}{L.mol^{-1}.s^{-1}}$，故 $P = 0.008$：不到百分之一的碰撞
具有正确的取向（NO 的 O 须与 \ce{O3} 的一个末端 O 相遇）。
\end{solution}

\begin{solution}{exo:b3:rate-theories:2}
$\Delta^{\ddagger}H^\circ = 75.0 - 2.48 = \qty{72.5}{kJ/mol}$（溶液中）；$75.0 - 4.96 = \qty{70.0}{kJ/mol}$
对双分子气相反应。
\end{solution}

\begin{solution}{exo:b3:rate-theories:3}
第尔斯–阿尔德反应：很大的负值（两个分子形成一个环状活化络合物）。
失去 \ce{N2}：正值（键松弛，自由度增加）。在
极性溶剂中生成离子：负值，正在形成的电荷使周围溶剂分子定向排列。
\end{solution}

\begin{solution}{exo:b3:rate-theories:4}
水：$8 \times 8.314 \times 298.15/(3 \times \num{0.890e-3}) = \qty{7.4e9}{L.mol^{-1}.s^{-1}}$；己烷
\qty{2.2e10}{L.mol^{-1}.s^{-1}}。己烷黏度只有水的三分之一，分子
扩散快三倍。
\end{solution}

\begin{solution}{exo:b3:rate-theories:5}
$\ln(k_2T_1/k_1T_2) = \ln(10 \times 298.15/318.15) = 2.238$; $\Delta^{\ddagger}H^\circ = 8.3145 \times 2.238/(1/298.15 -
1/318.15) = \qty{88.2}{kJ/mol}$. $\Delta^{\ddagger}S^\circ = R[\ln(kh/k_BT) + \Delta^{\ddagger}H^\circ/RT] = 8.3145 \times
(-36.365 + 35.596) = \qty{-6.4}{J.K^{-1}.mol^{-1}}$（在 \qty{298.15}{K} 下）。$E_a = R\ln10/(1/298.15 - 1/318.15) =
\qty{90.8}{kJ/mol} \approx \Delta^{\ddagger}H^\circ + RT$（在平均温度下）。
\end{solution}

\begin{solution}{exo:b3:rate-theories:6}
$\mu_{\mathrm{OH}} = 16/17$，$\mu_{\mathrm{OD}} = 32/18$，波数之比 $\sqrt{1.889} = 1.374$：
$\tilde\nu_{\mathrm{OD}} = \qty{2661}{cm^{-1}}$. $k_{\mathrm H}/k_{\mathrm D} = \exp(1.4388 \times 996/(2 \times 298)) = \eu^{2.40} = 11$.
\end{solution}

\begin{solution}{exo:b3:rate-theories:7}
最大值为 $\exp(1.4388 \times 808/(2 \times 298)) = \eu^{1.95} = 7.0$。数值 40 不可能仅由
零点能产生：氢原子穿过势垒隧穿，而氘
质量为其两倍，隧穿则少得多。
\end{solution}

\begin{solution}{exo:b3:rate-theories:8}
$\ce{S2O8^{2-}}$ 与 \ce{I-}：$z_{\mathrm A}z_{\mathrm B} = +2$，$\log(k/k_0) = 2 \times 0.51 \times 2 \times 0.10 = 0.20$：$k$
乘以 1.6。酯与 \ce{OH-}：一方为中性，无影响。
$\ce{[Co(NH3)5Br]^{2+}}$ 与 \ce{OH-}：$z_{\mathrm A}z_{\mathrm B} = -2$，$k$ 乘以 $10^{-0.20} = 0.63$。
\end{solution}

\begin{solution}{exo:b3:rate-theories:9}
较早出现，位于入口谷（过渡态类似反应物，据哈蒙德
假设）。逆反应吸热，其势垒在自身路径上靠后，
H--F 键中的振动能最能促进它。
\end{solution}

\begin{solution}{exo:b3:rate-theories:10}
$\Delta^{\ddagger}H^\circ = -Rb = \qty{61.9}{kJ/mol}$，标准误差为 $R \times 31 = \qty{0.26}{kJ/mol}$，
95\,\% 置信区间为 $\pm\qty{0.8}{kJ/mol}$。
\[
  \Delta^{\ddagger}S^\circ = R(16.528 - 23.760) = \qty{-60.1}{J.K^{-1}.mol^{-1}},
\]
标准误差为 $R \times 0.103 = \qty{0.86}{J.K^{-1}.mol^{-1}}$，95\,\% 置信区间为 $\pm\qty{2.7}{J.K^{-1}.mol^{-1}}$。
\end{solution}

\begin{solution}{exo:b3:rate-theories:11}
由 $\ln k = \ln(k_BT/h) + \Delta^{\ddagger}S^\circ/R - \Delta^{\ddagger}H^\circ/RT$，
\[
  \frac{\dd\ln k}{\dd T} = \frac1T + \frac{\Delta^{\ddagger}H^\circ}{RT^2}, \qquad
  E_a = RT^2\frac{\dd\ln k}{\dd T} = \Delta^{\ddagger}H^\circ + RT .
\] 于是 $A = k\eu^{E_a/RT} = (k_BT/h)\eu^{\Delta^{\ddagger}S^\circ/R}\eu^{-\Delta^{\ddagger}H^\circ/RT}\eu^{(\Delta^{\ddagger}H^\circ + RT)/RT} =
\eu\,(k_BT/h)\eu^{\Delta^{\ddagger}S^\circ/R}$。当 $\Delta^{\ddagger}S^\circ = 0$ 时：$A = 2.718 \times \num{6.21e12} = \qty{1.7e13}{s^{-1}}$，
与观测因子的数量级相当；高于此值则表明 $\Delta^{\ddagger}S^\circ$ 为正，
过渡态较为松散。
\end{solution}

\begin{solution}{exo:b3:rate-theories:12}
采用单位体积的分子配分函数：$q^{\mathrm T}/V = (2\pi mk_BT)^{3/2}/h^3$，适用于
A、B 及活化络合物（质量 $m_{\mathrm A} + m_{\mathrm B}$），对活化络合物则为 $q^{\mathrm R} = 8\pi^2Ik_BT/h^2$
（$\sigma = 1$）。平动之比为 $h^3/(2\pi\mu k_BT)^{3/2}$，于是
\[
  k = \frac{k_BT}{h}\cdot\frac{h^3}{(2\pi\mu k_BT)^{3/2}}\cdot\frac{8\pi^2\mu d^2k_BT}{h^2}\,\eu^{-E_0/RT}
    = \frac{8\pi^2d^2(k_BT)^{1/2}}{(2\pi)^{3/2}\mu^{1/2}}\eu^{-E_0/RT}
    = \pi d^2\sqrt{\frac{8k_BT}{\pi\mu}}\,\eu^{-E_0/RT},
\]
每对分子；乘以 $N_A$ 即得 \cref{thm:b3:rate-theories:collision}
其中 $\sigma = \pi d^2$。
\end{solution}

\begin{solution}{pb:b3:rate-theories:1}
\textbf{1.} $\mu_{\mathrm{CH}} = 12 \times 1.00783/13.00783 = \qty{0.9297}{u}$; $\mu_{\mathrm{CD}} = 12 \times 2.01410/14.01410 =
\qty{1.7246}{u}$.
\textbf{2.} $2917/\sqrt{1.7246/0.9297} = 2917/1.3620 = \qty{2142}{cm^{-1}}$，比
实测值 2109 高 1.5\,\%（非简谐性及 \ce{CD4} 中四个键的耦合）。
\textbf{3.} $\frac12 \times 2917 = \qty{1458.5}{cm^{-1}}$ 与 $\frac12 \times 2109 = \qty{1054.5}{cm^{-1}}$；差值
\qty{404}{cm^{-1}} $= \qty{4.83}{kJ/mol}$.
\textbf{4.} 它成为反应坐标，其零点能消失。
\textbf{5.} 对 D 而言，势垒高出 \qty{4.83}{kJ/mol}。
\textbf{6.} $\exp(1.43878 \times 404/310.15) = \eu^{1.874} = 6.5$.
\textbf{7.} $\exp(1.43878 \times 404/298.15) = 7.0$.
\textbf{8.} 在真实的过渡态中，伸缩振动零点能的一部分
保留于其他振动模式中，从而降低了此效应；仅由零点能
不可能超过此值（隧道效应可以）。
\textbf{9.} 很小：它们是次级效应，每个氘至多
约为 10\,\% 的量级。
\textbf{10.} 即 C--H 键在决速步中断裂，
其伸缩振动完全转化为反应坐标。
\textbf{11.} 一级消除反应：$t_{1/2} = \ln2/k_{\mathrm{total}}$，与之成反比。
\textbf{12.} $k_{\mathrm{total}}(\mathrm H) = 0.75\,k_{\mathrm{total}}(\mathrm H) + 0.25\,k_{\mathrm{total}}(\mathrm H)$，去甲基化及其他途径。
\textbf{13.} 除以 6.5。
\textbf{14.} $0.25 + 0.75/6.52 = 0.365$.
\textbf{15.} $5.0/0.365 = \qty{13.7}{h}$.
\textbf{16.} 2.7.
\textbf{17.} 6.5，即完整的同位素效应。
\textbf{18.} 剂量更小或次数更少，血液中的药物浓度更为平稳。
\textbf{19.} 以 $\ln(k/T)$ 对 $1/T$ 作图。
\textbf{20.} $\Delta^{\ddagger}H^\circ = 61.9 \pm \qty{0.3}{kJ/mol}$, $\Delta^{\ddagger}S^\circ = -60.1 \pm \qty{0.9}{J.K^{-1}.mol^{-1}}$.
\textbf{21.} $\Delta^{\ddagger}H^\circ = 66.8 \pm \qty{0.2}{kJ/mol}$, $\Delta^{\ddagger}S^\circ = -60.3 \pm \qty{0.7}{J.K^{-1}.mol^{-1}}$.
\textbf{22.} $\Delta\Delta^{\ddagger}H^\circ = 4.88 \pm \qty{0.33}{kJ/mol}$ ($\sqrt{0.26^2 + 0.21^2}$).
\textbf{23.} 是：4.83 落在一个标准误差之内，同位素效应
完全可由零点能解释，没有隧道效应的迹象。
\textbf{24.} 二者在误差范围内相等：过渡态对两种同位素异构体具有
相同的结构，如第一部分中所假设。
\textbf{25.} \textbf{$t_{1/2}(\mathrm D)/t_{1/2}(\mathrm H) \approx 2.7$，在 \qty{37}{\celsius} 时}：\qty{13.7}{h} 对 \qty{5.0}{h}。
\end{solution}
